\section{Conclusion}
\label{sec:conclusion}

This paper asked: why does a design object take one form rather than another? The answer proposed by \formlaere{} is that form is the position in morphospace where the shape grammar's response to the fitness landscape gradient; as perceived through the agents' cognitive cones; produced a locally stable poly-agentistic compromise.

The theory unifies three formal traditions: Stiny's shape algebra provides the syntax of form (what operations are possible), Hatchuel and Weil's C-K theory provides the epistemic logic (how knowledge expands through design), and Levin's multi-scale competence framework provides the dynamics (how agents navigate and why certain positions are stable). The synthesis produces a single generative equation: $\mathrm{Form} = SG(\nabla_{C(A)}\, L(c,t))$.

The theory introduces care; the dimensionality of the agent's cognitive cone; as the measure of design intelligence. An agent that represents more axes of the morphospace navigates a richer landscape and produces forms that withstand more pressures simultaneously. An agent that narrows its cone produces forms that fail along the axes it did not see. The failure is not accidental; it is systematic, and the form bears the signature of the blindness that produced it.

We tested the theory against 2,066 European chairs spanning 744 years. Twenty statistical analyses across 14 cross-validation subsets yield results consistent with the core predictions of Propositions 1--4; no postulate is falsified. Propositions 5--7 (agency, care, and the diagnostic conclusion) are formalized and illustrated but require direct empirical testing through the research programs proposed in Section~\ref{sec:agenda}.

We demonstrated the theory's reach beyond historical furniture through three worked examples: the cantilever chair race of 1925--1928, where independent convergence confirms landscape attractors; Neri Oxman's Silk Pavilion, where biological agents navigate the form grammar at cellular scale; and the Great Green Wall, where continental-scale ecological restoration demands care across incommensurable timescales.

The theory is falsifiable at every level. Each proposition generates testable predictions. Seven immediate research programs are proposed, each independently publishable and each capable of revising or refuting the theory.

Form shows how far care reached. The theory shows how to measure both.
